% titel: Blatt 6, A1: Elektronegativität vs. Elektronenaffinität

\subsection*{A) Unterschied + Vergleich der Paare}
\textbf{Elektronenaffinität (EA):} Energie, die frei wird, wenn ein isoliertes
Atom ein Elektron aufnimmt. \textbf{Elektronegativität (EN):} Fähigkeit eines
\textit{gebundenen} Atoms, Bindungselektronen an sich zu ziehen.
EN wächst nach rechts/oben im PSE, F ist am elektronegativsten.

Elektronegativere Atom je Paar:
\begin{center}
\begin{tabular}{c c c}
Paar & elektronegativer & Begründung \\
\hline
F / Cl      & \textbf{F}  & F steht rechts oben, F ist am elektronegativsten \\
F / O       & \textbf{F}  & F rechts von O in Periode 2 \\
B / Br      & \textbf{Br} & Br rechts von B und in Gruppe 17 hoch \\
O / Te      & \textbf{O}  & O rechts von Te und oben in Gruppe 16 \\
N / P       & \textbf{N}  & N rechts von P und oben in Gruppe 15 \\
\end{tabular}
\end{center}

\subsection*{B) Messbarkeit}
\textbf{Klar definiert und messbar:} die Elektronenaffinität (Energieänderung beim
isolierten Atom). \textbf{Nur vage/abschätzbar:} die Elektronegativität (relativer
Skalenwert, nicht direkt messbar).

\begin{ergebnis}
Elektronegativer: F, F, Br, O, N. Klar messbar: EA; nur abschätzbar: EN.
\end{ergebnis}

\begin{kontrolle}
Jeder bei B) genannte elektronegativere Partner liegt im PSE weiter rechts
bzw. oben als der andere --- konsistent mit rechts/oben-Wachstum der EN; die
Tabelle deckt sich mit den Pauling-Werten (z.\ B. F 3.98, Cl 3.16, O 3.44,
N 3.04, P 2.19, Br 2.96).
\end{kontrolle}

\begin{merksatz}
EA = freies Atom, messbar; EN = gebundenes Atom, nur Skalenwert. Elektronegativ
heisst nach rechts/oben; F führt die Skala an.
\end{merksatz}
